By Proposition \claimref{singularity.stretched-singularity-count}, which holds under (C1) (\claimref{singularity.c1-curvature-cutoff}: the boundary mode space is supported on the slice $r = r_c = (2M)^{1/3}$) and (C2) (\claimref{singularity.c2-cell-density}: $\nu$ fermionic modes per Planck cell of that slice), the modes accrue at $dN/dt = 8\pi\nu M\sqrt{1 - r_c/2M}$. The number of Planck cells accumulated per unit Killing time is the slice volume of Proposition \claimref{scrambling.slice-volume-exactly-per-unit-killing} divided by $\lP^3 = 1$, i.e.\ $8\pi M$ at leading order, the leading term of the same factor; multiplying by $\nu$ and by the span $\Delta t$ gives $N$. The second display substitutes $\TH = 1/8\pi M$, which Proposition \claimref{singularity.stretched-singularity-count} takes from modular descent under (R) (\claimref{singularity.r-regularity-fixed-locus}); the count $N$ itself uses only (C1) and (C2). The corank factor in $\nu$ is fixed at one by (C2); what remains in $\nu$ is the Planck-cell calibration and the bit or nat convention (Remark \ref{rem:nu}), and the proportionality $dN/dt \propto 1/\TH$ does not depend on it. The placement of the modes on the slice, and their density per cell, are the two postulates and nothing else: in particular the tower's Witt-index capacity, which counts mode pairs per rung, supplies neither, and the modes counted here are the ladder pairs of (C2), one per point of the blown-up singular boundary, not rungs of the tower.
